\begin{frame}[fragile]{실습 1: 죄수의 딜레마 자기학습}
  두 에이전트가 같은 보상 행렬을 놓고, \textbf{둘 다 Week 6의
  $\varepsilon$-greedy Q-learning으로 배우며 서로를 상대로
  플레이(self-play)} --- 알고리즘은 각자가 자기 $Q$값을
  갱신할 뿐(상대도 학습 중이라는 사실은 \textbf{어디에도
  명시되지 않음}):
\begin{lstlisting}
import numpy as np
# 행=에이전트1 행동, 열=에이전트2 행동 (0=C 침묵, 1=D 자백)
PD = np.array([[3.0, 0.0],   # (C,C)=3, (C,D)=0
               [5.0, 1.0]])  # (D,C)=5, (D,D)=1

def eps_greedy(Q, eps, rng):
    if rng.random() < eps:
        return int(rng.integers(len(Q)))
    return int(np.argmax(Q))

def self_play(games, payoff, alpha=0.5, eps_start=1.0, eps_end=0.02,
              tau=500.0, seed=0):
    rng = np.random.default_rng(seed)
    Q1, Q2 = np.zeros(2), np.zeros(2)
    rewards1 = []
    for g in range(games):
        eps = eps_end + (eps_start - eps_end) * np.exp(-g / tau)  # 탐험 감소
        a1 = eps_greedy(Q1, eps, rng)
        a2 = eps_greedy(Q2, eps, rng)
        r1 = payoff[a1, a2]
        r2 = payoff[a2, a1]               # 게임이 대칭
        Q1[a1] += alpha * (r1 - Q1[a1])   # Week 6의 TD 갱신 그대로
        Q2[a2] += alpha * (r2 - Q2[a2])
        rewards1.append(r1)
    return np.array(rewards1), Q1, Q2

rewards, Q1, Q2 = self_play(20000, PD)
\end{lstlisting}
\end{frame}
